\documentclass[10pt,a4paper]{amsart}
\usepackage{amsmath,amssymb,amsthm,graphicx}
% Retained minimal preamble for the Stacks Project chapter more-algebra.tex (More on Algebra).
% The chapter file pinned-original-source.tex carries only \input{preamble} on its first line, so no
% preamble text lives inside the chapter itself.  The Stacks Project uses one single shared preamble for
% every chapter; the snapshot stages it identically as sources/stacks-chow/preamble.tex and
% sources/stacks-etale-cohomology/preamble.tex.  The block below is copied mechanically, line for line,
% from sources/stacks-chow/preamble.tex lines 14-18 (the Xy-pic setup) and 154-217 (the native theorem
% declarations together with the mathematical macros).  Nothing was added, removed or redefined; the
% wrapper class and the amsmath/amssymb/amsthm/graphicx base packages are supplied by the assembler.
% Chapter: more-algebra.tex -- More on Algebra.
% For commutative diagrams we use Xy-pic
\usepackage[all]{xy}

% We use 2cell for 2-commutative diagrams.
\xyoption{2cell}

% Theorem environments.
%
\theoremstyle{plain}
\newtheorem{theorem}[subsection]{Theorem}
\newtheorem{proposition}[subsection]{Proposition}
\newtheorem{lemma}[subsection]{Lemma}

\theoremstyle{definition}
\newtheorem{definition}[subsection]{Definition}
\newtheorem{example}[subsection]{Example}
\newtheorem{exercise}[subsection]{Exercise}
\newtheorem{situation}[subsection]{Situation}

\theoremstyle{remark}
\newtheorem{remark}[subsection]{Remark}
\newtheorem{remarks}[subsection]{Remarks}

\numberwithin{equation}{subsection}

% Macros
%
\def\lim{\mathop{\mathrm{lim}}\nolimits}
\def\colim{\mathop{\mathrm{colim}}\nolimits}
\def\Spec{\mathop{\mathrm{Spec}}}
\def\Hom{\mathop{\mathrm{Hom}}\nolimits}
\def\Ext{\mathop{\mathrm{Ext}}\nolimits}
\def\SheafHom{\mathop{\mathcal{H}\!\mathit{om}}\nolimits}
\def\SheafExt{\mathop{\mathcal{E}\!\mathit{xt}}\nolimits}
\def\Sch{\mathit{Sch}}
\def\Mor{\mathop{\mathrm{Mor}}\nolimits}
\def\Ob{\mathop{\mathrm{Ob}}\nolimits}
\def\Sh{\mathop{\mathit{Sh}}\nolimits}
\def\NL{\mathop{N\!L}\nolimits}
\def\CH{\mathop{\mathrm{CH}}\nolimits}
\def\proetale{{pro\text{-}\acute{e}tale}}
\def\etale{{\acute{e}tale}}
\def\QCoh{\mathit{QCoh}}
\def\Ker{\mathop{\mathrm{Ker}}}
\def\Im{\mathop{\mathrm{Im}}}
\def\Coker{\mathop{\mathrm{Coker}}}
\def\Coim{\mathop{\mathrm{Coim}}}

% Boxtimes
%
\DeclareMathSymbol{\boxtimes}{\mathbin}{AMSa}{"02}

%
% Macros for moduli stacks/spaces
%
\def\QCohstack{\mathcal{QC}\!\mathit{oh}}
\def\Cohstack{\mathcal{C}\!\mathit{oh}}
\def\Spacesstack{\mathcal{S}\!\mathit{paces}}
\def\Quotfunctor{\mathrm{Quot}}
\def\Hilbfunctor{\mathrm{Hilb}}
\def\Curvesstack{\mathcal{C}\!\mathit{urves}}
\def\Polarizedstack{\mathcal{P}\!\mathit{olarized}}
\def\Complexesstack{\mathcal{C}\!\mathit{omplexes}}
% \Pic is the operator that assigns to X its picard group, usage \Pic(X)
% \Picardstack_{X/B} denotes the Picard stack of X over B
% \Picardfunctor_{X/B} denotes the Picard functor of X over B
\def\Pic{\mathop{\mathrm{Pic}}\nolimits}
\def\Picardstack{\mathcal{P}\!\mathit{ic}}
\def\Picardfunctor{\mathrm{Pic}}
\def\Deformationcategory{\mathcal{D}\!\mathit{ef}}

% --- appended 2026-10-12 (Chief): the project own definitions of the reference/slogan/history
% display environments, copied mechanically from the shared Stacks preamble lines 8-12
% (\usepackage{verbatim} plus the three \newenvironment lines, which the chapter preamble keeps
% verbatim from the project preamble). Nothing added, removed or redefined: these environments
% carry no mathematical content and the project itself typesets them as comments. Without these
% three lines a counted statement whose own text uses them fails with Environment slogan undefined.
\usepackage{verbatim}
\newenvironment{reference}{\comment}{\endcomment}
\newenvironment{slogan}{\comment}{\endcomment}
\newenvironment{history}{\comment}{\endcomment}
\usepackage{cleveref}
\crefname{originalequation}{original equation}{original equations}
\Crefname{originalequation}{Original equation}{Original equations}
\makeatletter
\newlabel{equation-base-change-RHom}{{0E1X}{source}{}{sourcepointer-1}{}}
\newlabel{equation-base-change-RHom@cref}{{[originalequation][0][]0E1X}{[1][0][]source}{}{}{}}
\makeatother
\crefname{originallemma}{original lemma}{original lemmas}
\Crefname{originallemma}{Original lemma}{Original lemmas}
\makeatletter
\newlabel{lemma-internal-hom-diagonal-better}{{0BYN}{source}{}{sourcepointer-2}{}}
\newlabel{lemma-internal-hom-diagonal-better@cref}{{[originallemma][0][]0BYN}{[1][0][]source}{}{}{}}
\makeatother
\crefname{originallemma}{original lemma}{original lemmas}
\Crefname{originallemma}{Original lemma}{Original lemmas}
\makeatletter
\newlabel{lemma-upgrade-adjoint-tensor-RHom}{{0E1W}{source}{}{sourcepointer-3}{}}
\newlabel{lemma-upgrade-adjoint-tensor-RHom@cref}{{[originallemma][0][]0E1W}{[1][0][]source}{}{}{}}
\makeatother
\crefname{originallemma}{original lemma}{original lemmas}
\Crefname{originallemma}{Original lemma}{Original lemmas}
\makeatletter
\newlabel{lemma-internal-hom-evaluate-tensor-isomorphism}{{0ATK}{source}{}{sourcepointer-4}{}}
\newlabel{lemma-internal-hom-evaluate-tensor-isomorphism@cref}{{[originallemma][0][]0ATK}{[1][0][]source}{}{}{}}
\makeatother
\title{Base change compatibility of RHom under the four standard hypotheses}
\author{The Stacks Project Authors}
\date{stacks-project:more-algebra.tex:local-snapshot-20261012; GNU FDL 1.2 or later}
\begin{document}
\maketitle
\setcounter{subsection}{1}
\noindent\textit{Editorial note.} All delivered mathematical text is the author's own, copied from the pinned Stacks Project chapter file more-algebra.tex; no statement or proof step was rewritten, reordered or removed. The counted span is the author environment and its own proof environment. References inside the span are handed to the assembler as navigation stubs that carry the canonical tag of the cited tag only; whether each of them is load-bearing is recorded on the stub itself.\par\medskip
\noindent\textit{Allowed standard prerequisites (editorial).}\begin{itemize}
\item The base change map for RHom is characterized by the adjunction formula relating it to the tensor-internal-hom diagonal map. invoked at proof line 28851. Matches the restricted map with the diagonal map of the next lemma.
\item The tensor-internal-hom diagonal map K tensor RHom(M, L) -> RHom(M, K tensor L) is an isomorphism when K is perfect, R' is perfect over R, the map is flat with K pseudo-coherent and M bounded below, or R' has finite tor dimension with K pseudo-coherent and M bounded below. invoked at proof lines 28852-28853. Supplies the isomorphism in each of the four listed cases.
\end{itemize}\medskip
\ifdefined\hypertarget\hypertarget{sourcepointer-1}{}\fi
\noindent\textit{Source reference.} Original equation 0E1X: sources/stacks-more-algebra/pinned-original-source.tex, source lines 28801--28801. Referenced at statement line 28822 as the map whose isomorphy is claimed; the map is the subject inside the counted span.\par\medskip
\ifdefined\hypertarget\hypertarget{sourcepointer-2}{}\fi
\noindent\textit{Source reference.} Original lemma 0BYN: sources/stacks-more-algebra/pinned-original-source.tex, source lines 18036--18036. Cited at proof line 28847 to identify the restricted map with the diagonal map; the conclusion is used inside the counted span.\par\medskip
\ifdefined\hypertarget\hypertarget{sourcepointer-3}{}\fi
\noindent\textit{Source reference.} Original lemma 0E1W: sources/stacks-more-algebra/pinned-original-source.tex, source lines 28763--28763. Cited at proof line 28851 to match the left and right hand sides of the restricted map; the conclusion is used inside the counted span.\par\medskip
\ifdefined\hypertarget\hypertarget{sourcepointer-4}{}\fi
\noindent\textit{Source reference.} Original lemma 0ATK: sources/stacks-more-algebra/pinned-original-source.tex, source lines 28653--28653. Cited at proof line 28852 to conclude the isomorphism in each listed case; the conclusion is used inside the counted span.\par\medskip
\begin{lemma}
\label{lemma-base-change-RHom}
Let $R \to R'$ be a ring map. Let $K, M \in D(R)$. The map
(\ref{equation-base-change-RHom})
$$
R\Hom_R(K, M) \otimes_R^\mathbf{L} R'
\longrightarrow
R\Hom_{R'}(K \otimes_R^\mathbf{L} R', M \otimes_R^\mathbf{L} R')
$$
is an isomorphism in $D(R')$ in the following cases
\begin{enumerate}
\item $K$ is perfect,
\item $R'$ is perfect as an $R$-module,
\item $R \to R'$ is flat, $K$ is pseudo-coherent, and $M \in D^{+}(R)$, or
\item $R'$ has finite tor dimension as an $R$-module,
$K$ is pseudo-coherent, and $M \in D^{+}(R)$
\end{enumerate}
\end{lemma}

\begin{proof}
We may check the map is an isomorphism after applying the
restriction functor $D(R') \to D(R)$. After applying this
functor our map becomes the map
$$
R\Hom_R(K, L) \otimes_R^\mathbf{L} R'
\longrightarrow
R\Hom_R(K, L \otimes_R^\mathbf{L} R')
$$
of Lemma \ref{lemma-internal-hom-diagonal-better}.
See discussion above the lemma to match the left and right hand sides;
in particular, this uses Lemma \ref{lemma-upgrade-adjoint-tensor-RHom}.
Thus we conclude by
Lemma \ref{lemma-internal-hom-evaluate-tensor-isomorphism}.
\end{proof}

\end{document}
